\documentclass[10pt,a4paper]{article}

% ----- encoding & layout -----
\usepackage[utf8]{inputenc}
\usepackage[T1]{fontenc}
\usepackage{lmodern}
\usepackage[a4paper,margin=2cm]{geometry}
\usepackage{microtype}
\usepackage{parskip}

% ----- tables -----
\usepackage{array}
\usepackage{longtable}
\usepackage{booktabs}
\usepackage[table]{xcolor}
\usepackage{ragged2e}
\usepackage{pdflscape}

% ----- misc -----
\usepackage{pifont}
\usepackage{enumitem}
\usepackage{fancyhdr}
\usepackage[colorlinks=true,linkcolor=blue!50!black,urlcolor=blue!45!black,citecolor=blue!50!black]{hyperref}

% ----- header / footer -----
\setlength{\headheight}{13pt}
\pagestyle{fancy}
\fancyhf{}
\lhead{\small Barrel / Cylinder Detection --- Methods Survey}
\rhead{\small p.\,\thepage}
\lfoot{\small \texttt{researchwrite/} --- living document}
\rfoot{\small compiled 2026-05-30}
\renewcommand{\headrulewidth}{0.4pt}
\renewcommand{\footrulewidth}{0.4pt}

% ----- rating colours (light fills, black text) -----
\definecolor{cVH}{HTML}{81C784} % very high  - green
\definecolor{cH}{HTML}{C5E1A5}  % high       - light green
\definecolor{cMod}{HTML}{FFE082} % moderate  - amber
\definecolor{cLM}{HTML}{FFCC80}  % low-mod   - orange
\definecolor{cLow}{HTML}{EF9A9A} % low       - red

% ----- star rating 1..5 -----
\newcommand{\stars}[1]{%
  \ifcase#1\or
  \ding{72}\ding{73}\ding{73}\ding{73}\ding{73}\or
  \ding{72}\ding{72}\ding{73}\ding{73}\ding{73}\or
  \ding{72}\ding{72}\ding{72}\ding{73}\ding{73}\or
  \ding{72}\ding{72}\ding{72}\ding{72}\ding{73}\or
  \ding{72}\ding{72}\ding{72}\ding{72}\ding{72}\fi}

% ragged-right fixed-width column
\newcolumntype{R}[1]{>{\RaggedRight\arraybackslash}p{#1}}

\begin{document}

% =================== TITLE ===================
\begin{center}
  {\LARGE\bfseries Barrel / Cylinder Detection in 3D Point Clouds}\\[2pt]
  {\large A Survey of Methods and their Suitability for Occluded Detection}\\[6pt]
  {\normalsize Master's-thesis literature backbone \quad$\cdot$\quad Bharath \quad$\cdot$\quad 2026-05-30}\\[2pt]
  {\small Living document --- regenerated as the method list and references evolve.}
\end{center}
\vspace{2pt}
\hrule
\vspace{8pt}

% =================== 1. PURPOSE ===================
\section{Purpose and scope}

This document is the literature backbone for the thesis comparison of barrel/cylinder
detection methods in 3D point clouds. The aim is \textbf{empirical and comparative}: to
determine \emph{which methods detect barrels most accurately} --- particularly under
occlusion and partial burial --- and to \emph{characterise the strengths, weaknesses, and
operating regimes} of each method family.

It is deliberately \textbf{not} committed to proving a single predetermined hypothesis.
Any trade-offs we observe (for example, how precision and recall move as occlusion
increases, or how camera+LiDAR fusion compares with LiDAR-only methods) are
\emph{findings to be measured}, not assumptions to be defended.

The comparison is made fair by a shared data / I/O contract: every method reads the same
clouds and emits the same prediction schema, and a single evaluator scores them
(precision / recall / F1, radius and centre error, axis-angle error, runtime, and
annotation / training cost kept as a column so geometric and learned methods compare
fairly). Ground-truthed data --- including the \emph{true per-barrel occlusion percentage}
--- will come from NVIDIA Isaac Sim, with validation on a small set of real Asus Xtion
scans. The current baseline is the 3DTK randomized-Hough \texttt{detectCylinder}.

\textbf{Reading the survey.} Methods already named in the project plan are marked with a
dagger~($\dagger$); the rest were surfaced by a literature search on 2026-05-30. Every
method maps to one or more entries in the bibliography (\S6) via the short key in the
\emph{Key refs} column.

% =================== 2. RUBRIC ===================
\section{How to read the occlusion-suitability rating}

Each method carries a 1--5 rating for \textbf{suitability to occluded / partial-view
barrel detection} --- the central difficulty of this thesis. The rating is an informed
prior from the literature whose only job is to \emph{prioritise experiments}; the thesis
will replace it with measured numbers. It weighs five criteria:

\begin{enumerate}[leftmargin=1.6em,itemsep=1pt,topsep=2pt]
  \item \textbf{Partial-evidence detection} --- can it fire from a partial arc or sparse returns?
  \item \textbf{Robustness to fragmentation} --- does missing surface break its core mechanism (region growth, clustering, connectivity)?
  \item \textbf{Learnability of occlusion} --- can occluded cases be learned from data, and is such data affordable here?
  \item \textbf{Second-modality / completion support} --- does a camera or a completion stage recover the object when LiDAR is occluded?
  \item \textbf{Localisation under clutter} --- can it separate adjacent / overlapping barrels?
\end{enumerate}

\noindent\textbf{Scale.}\quad
\colorbox{cVH}{\textbf{5} Very High} purpose-built for occlusion \;\textbar\;
\colorbox{cH}{\textbf{4} High} handles it well (often data-dependent) \;\textbar\;
\colorbox{cMod}{\textbf{3} Moderate} partial views OK, degrades with heavy occlusion/clutter \;\textbar\;
\colorbox{cLM}{\textbf{2} Low--Mod} fragile under occlusion \;\textbar\;
\colorbox{cLow}{\textbf{1} Low} unsuited.

% =================== 3. THE TABLE ===================
\begin{landscape}
\section{Methods comparison table}
\footnotesize
\setlength{\extrarowheight}{2.5pt}
\renewcommand{\arraystretch}{1.05}

\begin{longtable}{@{}R{0.5cm} R{3.3cm} R{6.7cm} R{2.2cm} R{7.0cm} R{2.5cm}@{}}
\rowcolor{gray!35}
\textbf{\#} & \textbf{Method} & \textbf{Core idea} & \textbf{Occlusion fit} & \textbf{Why (behaviour under occlusion / clutter)} & \textbf{Key refs}\\
\endhead

% ---- Group A ----
\rowcolor{black!8}\multicolumn{6}{@{}l}{\textbf{A.\quad LiDAR-only --- classical / geometric}}\\
1 & Euclidean clustering + per-cluster fit~$\dagger$ &
Remove ground/planes, cluster the rest (DBSCAN/Euclidean), fit one cylinder per cluster. &
\cellcolor{cLM}\stars{2}\newline Low--Mod &
Occlusion and adjacency split or merge clusters; touching barrels collapse into one. Fast and simple, but the proposer is the weak link. \emph{(Project's current proposer.)} &
project pipeline\\

2 & Region growing + normal curvature~$\dagger$ &
Grow surface regions under a smoothness / curvature constraint, then fit. &
\cellcolor{cLM}\stars{2}\newline Low--Mod &
Needs a contiguous smooth surface; occlusion fragments the shell into disconnected patches, giving broken / under-segmented regions. &
Rabbani'06, Moradi'21\\

3 & Randomized 2-step Hough --- 3DTK \texttt{detectCylinder} \emph{(baseline)}~$\dagger$ &
2D Hough on the Gaussian sphere of normals for the axis direction, then 3D Hough for centre + radius. &
\cellcolor{cMod}\stars{3}\newline Moderate &
A partial arc still votes and axis voting tolerates partial views; but heavy occlusion thins the accumulator and clutter creates spurious peaks. Lock radius to the drum band; tune accumulator resolution + sampling. &
Rabbani'05, Borrmann'11\\

4 & RANSAC cylinder fit (PCL)~$\dagger$ &
Sample minimal point+normal sets, hypothesise a cylinder, score inliers; PCA for axis. &
\cellcolor{cMod}\stars{3}\newline Moderate &
Tolerates outliers and fits from a partial arc, but inlier thresholds are noise-sensitive and clutter spawns competing hypotheses; needs a proposer to localise. &
PCL; Nurunnabi'19\\

5 & Efficient RANSAC --- multi-primitive &
Accelerated RANSAC detecting planes / spheres / cylinders / cones / tori at once. &
\cellcolor{cMod}\stars{3}\newline Moderate &
Robust to heavy noise/outliers and scales well; partial cylinders detectable, but over-detects primitives, so it needs post-filtering to the barrel prior. &
Schnabel'07\\

6 & Robust least-squares cylinder fit \emph{(fit stage)} &
Faithful / robust (M-estimator, robust PCA) least-squares fit of cylinder parameters to inliers. &
\cellcolor{cMod}\stars{3}\newline Moderate &
Degrades gracefully on partial arcs and resists outliers --- but it is a \emph{fit}, not a detector; quality is bounded by the proposer that hands it points. &
Luk\'acs'98, Nurunnabi'19\\

7 & Connectivity / projection-based detection &
Project points along hemisphere directions; detect circular projections via connected components. Little tuning, arbitrary orientation. &
\cellcolor{cMod}\stars{3}\newline Moderate &
Noise-robust and parameter-light; circular-projection evidence survives partial views, but connected-component reasoning can break under fragmentation. &
Ara\'ujo'20\\

% ---- Group B ----
\rowcolor{black!8}\multicolumn{6}{@{}l}{\textbf{B.\quad LiDAR-only --- learned}}\\
8 & PointNet++ (segment $\to$ fit)~$\dagger$ &
Hierarchical point-set network segments barrel points; metric fit recovers pose/radius. &
\cellcolor{cH}\stars{4}\newline High &
Learns partial-shape cues, so it can flag occluded barrels \emph{if trained on occluded examples} --- exactly BarrelNet's recipe. Cost: annotation + sim-to-real gap. &
Qi'17; cf.\ Yan'24\\

9 & VoteNet --- deep Hough voting &
Seed points vote for object centres; cluster votes into proposals. &
\cellcolor{cH}\stars{4}\newline High &
Centre-voting is built for partial / occluded objects --- visible surface points still vote toward the hidden centroid. Needs 3D-box training labels. &
Qi'19\\

10 & Voxel / pillar detectors --- PointPillars, SECOND, PV-RCNN~$\dagger$ &
Voxel/pillar encoder + sparse-conv backbone + detection head; PV-RCNN adds point-voxel refinement. &
\cellcolor{cH}\stars{4}\newline High &
Trained end-to-end on occlusion-rich benchmarks; learn to detect from partial evidence. Strong accuracy; cost = annotated data + compute. &
Lang'19, Shi'20\\

11 & CenterPoint --- center-based detection~$\dagger$ &
Detect object centres as BEV heatmap peaks; regress box / size / orientation. &
\cellcolor{cH}\stars{4}\newline High &
The centre representation is robust to partial observation and crowding; SOTA-class on driving data. Same data / compute caveat. &
Yin'21\\

12 & BtcDet --- ``Behind the Curtain'' (occlusion-aware) &
Learn shape priors, estimate occupancy in occluded / missed regions, complete shapes, then detect. &
\cellcolor{cVH}\stars{5}\newline Very High &
Purpose-built for occlusion --- explicitly reasons about what is hidden. Highest relevance to the occlusion question; higher engineering complexity. &
Xu'22\\

13 & Learned primitive fitting --- SPFN / Point2Cyl / PrimiTect &
Predict per-point primitive membership/type/normals (SPFN), extrusion-cylinder decomposition (Point2Cyl), or continuous Hough voting (PrimiTect). &
\cellcolor{cMod}\stars{3}\newline Moderate &
Beat RANSAC on clean / CAD data and give principled cylinder params; partial-view + sim-to-real generalisation for field barrels is unproven. Point2Cyl targets CAD extrusions. &
Li'19, Uy'22, Sommer'20\\

14 & Shape completion $\to$ fit --- PCN \emph{(pre-stage)} &
Complete the partial cloud (encoder-decoder), then fit / detect on the completed shell. &
\cellcolor{cH}\stars{4}\newline High &
Directly attacks occlusion by hallucinating the missing shell before fitting; adds a learned stage with its own failure modes and training needs. &
Yuan'18\\

% ---- Group C ----
\rowcolor{black!8}\multicolumn{6}{@{}l}{\textbf{C.\quad Camera + LiDAR fusion}}\\
15 & Camera-first frustum --- YOLO $\to$ frustum $\to$ fit; Frustum-PointNets / -ConvNet~$\dagger$ &
2D image detector proposes a region; back-project to a 3D frustum; segment + fit inside it. &
\cellcolor{cH}\stars{4}\newline High &
The camera still sees a barrel when its LiDAR return is sparse / occluded $\Rightarrow$ strong recall under occlusion. Failure modes: object must be visible in-image, and extrinsic-calibration error degrades the frustum. &
Qi'18, Wang'19\\

16 & LiDAR-first $\to$ image verification~$\dagger$ &
Generate 3D proposals from LiDAR, project to the image, verify / score with the 2D detector (decision-level fusion). &
\cellcolor{cH}\stars{4}\newline High &
Image confirmation suppresses false positives and rescues low-confidence detections; but the proposal still needs enough LiDAR points on the occluded object. &
Pang'20 (CLOCs)\\

17 & Point decoration --- PointPainting / PointAugmenting~$\dagger$ &
Paint each LiDAR point with image semantics / CNN features, then run any LiDAR detector. &
\cellcolor{cH}\stars{4}\newline High &
Semantic painting separates barrels from clutter and helps thin returns; effectiveness still bounded by having points on the object and by calibration. &
Vora'20\\

18 & Deep multimodal fusion --- TransFusion / BEVFusion / EPNet / MVX-Net &
Feature-level fusion of camera + LiDAR (transformer queries / unified BEV). &
\cellcolor{cVH}\stars{5}\newline Very High &
SOTA on occluded driving benchmarks and comparatively calibration-robust (TransFusion's soft association); cost = training data, compute, integration complexity. &
Bai'22, Liu'23\\

\end{longtable}

\vspace{2pt}
{\footnotesize $\dagger$ = method already named in the project plan (\texttt{CLAUDE.md} / memory). \quad
\emph{(baseline)} = current working baseline. \quad Rating scale defined in \S2.}
\end{landscape}

% =================== 4. PRIOR WORK ===================
\section{Closest prior work: occlusion and barrels specifically}

A targeted search surfaced work sitting almost exactly on this thesis:

\begin{itemize}[leftmargin=1.4em,itemsep=3pt,topsep=2pt]
  \item \textbf{BarrelNet} (Yan et al., OCEANS 2024) is the nearest neighbour: a
  PointNet-style network trained on \emph{synthetically generated} cylinder point clouds
  that explicitly model camera-viewpoint occlusion \emph{and} burial in the seafloor. It
  regresses 6-DOF pose, radius, and a \emph{burial fraction}, and \emph{beats a classical
  least-squares cylinder fit} on synthetic tests, with qualitative ROV-footage
  generalisation at a Southern-California DDT-barrel dump site. This validates the core
  thesis bet --- synthetic occluded / buried training plus comparison against a geometric
  baseline --- and is the single most important paper to position against.

  \item \textbf{Behind the Curtain / BtcDet} (Xu et al., AAAI 2022) makes occlusion a
  first-class citizen for generic LiDAR detection: it learns shape priors and predicts
  occupancy in occluded regions \emph{before} detecting. It is the strongest template for
  an occlusion-aware learned detector here.

  \item \textbf{Shape completion} (PCN, Yuan et al., 3DV 2018) is an orthogonal lever:
  complete the occluded shell, \emph{then} fit --- a drop-in pre-stage for any geometric
  method.

  \item \textbf{Occlusion analysis} --- OccAM's Laser (Schinagl et al., CVPR 2022)
  produces occlusion-based attribution maps for 3D detectors, useful for \emph{explaining}
  why a method's recall falls as occlusion rises.

  \item \textbf{Buried-object context beyond LiDAR} --- integrated geophysical surveys
  (magnetometry / electromagnetic induction) detect buried steel drums where optical /
  LiDAR sensing cannot. Out of scope as a detection method, but relevant when framing the
  buried-barrel problem and its sensing limits.
\end{itemize}

\noindent\textbf{Takeaway.} Occlusion-aware learned detection and synthetic occluded
training are established, and in the one published barrel-specific study a learned method
beats classical fitting. The thesis should therefore treat a learned occlusion-aware
method (a BarrelNet-style PointNet++, or BtcDet) as a \emph{serious contender} against the
3DTK Hough baseline, not an afterthought --- while the geometric methods remain the
zero-training reference point.

\subsection{Quantitative motivation: occlusion sensitivity and LiDAR-only reliability
limits}

Beyond the barrel-specific work above, two broader, quantitatively documented findings
from the general 3D-detection literature motivate why this thesis treats
occlusion-robustness and sensor fusion as central questions rather than incidental ones:

\begin{itemize}[leftmargin=1.4em,itemsep=3pt,topsep=2pt]
  \item \textbf{Detection accuracy degrades sharply and predictably with occlusion
  severity.} On KITTI, whose Easy/Moderate/Hard splits are explicitly defined by
  occlusion/truncation level, PointRCNN \textbf{[Shi'19]} drops from 88.88\% to 78.63\% to
  77.38\% car AP$_{3D}$@0.7 IoU (val set) across the three tiers --- an
  $\sim$11.5-point gap attributable to visibility alone, a pattern that recurs across
  nearly every LiDAR detector subsequently benchmarked on KITTI. A 2025 controlled-
  occlusion study on nuScenes \textbf{[Kumar'25]} makes the same point with directly
  injected occlusion: LiDAR-only mAP collapses from 64.7\% to 34.1\% ($-$47.3\%) under
  heavy sensor occlusion, and camera-only mAP drops from 35.6\% to 20.9\% ($-$41.3\%)
  under moderate occlusion.

  \item \textbf{LiDAR-only perception is measurably less reliable than LiDAR+camera
  fusion, including before occlusion is even considered.} In the same study
  \textbf{[Kumar'25]}, fused BEVFusion outperforms LiDAR-only on clean, non-occluded data
  (68.5\% vs 64.7\% mAP) and is far more robust when one sensor degrades: occluding the
  camera leg costs fusion only 4.1 points (68.5\%$\to$65.7\%), versus the 47.3-point
  LiDAR-only collapse above. The caveat that matters for a mission-critical framing:
  fusion still loses 26.8 points (68.5\%$\to$50.1\%) when the LiDAR leg itself is
  occluded --- current fusion architectures remain disproportionately reliant on LiDAR
  and \emph{mitigate}, not \emph{eliminate}, the single-sensor occlusion failure mode. A
  broader survey of multi-sensor fusion for autonomous driving \textbf{[Wang'25]} reaches
  the same qualitative conclusion: LiDAR alone is limited by occlusion and cost, cameras
  alone fail in low light / adverse weather and lack depth, and only fusion is positioned
  to cover both gaps.
\end{itemize}

\noindent\textbf{Implication for this thesis.} These are general findings from the wider
3D-detection literature, not yet measurements from this project's own data --- but they
are why \S1's open question (how camera+LiDAR fusion compares with LiDAR-only,
especially under occlusion) is not rhetorical: the literature already shows occlusion
substantially degrades LiDAR-only detection and that fusion only partially compensates.
The thesis's own occlusion sweep (\S5) is designed to produce the barrel-specific
equivalent of these numbers.

% =================== 5. SHORTLIST ===================
\section{Suggested shortlist for the comparison}

Balancing relevance, occlusion-suitability, and implementation cost, a defensible first
comparison set is:

\begin{itemize}[leftmargin=1.4em,itemsep=2pt,topsep=2pt]
  \item \textbf{Geometric baselines:} 3DTK randomized Hough \emph{(already running)} and a
  RANSAC / efficient-RANSAC cylinder fit --- cheap, no training, the reference point.
  \item \textbf{Learned LiDAR:} a PointNet++ / BarrelNet-style detector trained on
  Isaac-Sim occluded / buried barrels; optionally BtcDet as the occlusion specialist.
  \item \textbf{Fusion:} one camera-first frustum pipeline (YOLO $\to$ frustum $\to$ fit)
  to probe recall under occlusion, with a calibration-perturbation analysis as the fusion
  stress test.
\end{itemize}

\noindent This spans all three families, includes the occlusion specialist, and keeps the
geometric reference --- enough to answer \emph{which method works best, and why, per
occlusion regime}.

% =================== 5b. ML BY SUPERVISION ===================
\section{Machine-learning methods by supervision: what fits \emph{this} data}
\label{sec:mlsupervision}

\emph{(Added 2026-06-30.)} Sections~3--5 survey methods by family and occlusion-suitability.
This section re-cuts the \emph{learned} methods by the question that actually decides what
the thesis can build: \textbf{can we train and run it on the data we have?} It splits the
candidates into \textbf{supervised} and \textbf{unsupervised / self-supervised}, and maps
each onto the project's data reality and proposer$+$fit architecture.

\subsection{Three data constraints that decide viability}
\begin{enumerate}[leftmargin=1.6em,itemsep=2pt,topsep=2pt]
  \item \textbf{No real labels yet} (one annotated drum so far) and only a handful of real
  scans $\Rightarrow$ supervised-from-scratch on real data is out. The realistic path is
  \emph{train on synthetic} (\texttt{common/synth\_cylinder.py} now, NVIDIA Isaac Sim later,
  both with free per-point / instance / 6-DOF ground truth) and \emph{transfer to real} ---
  which makes the \emph{sim-to-real gap} a quantity to \emph{measure}, not avoid
  [Saleh'20].
  \item \textbf{Single, static, overhead viewpoint} (mostly top caps + short walls + radial
  scan-shadow). This excludes the large branch of unsupervised 3D detection that mines
  objects from \emph{motion across frames} (LISO~[Baur'24] and its relatives need sequential
  LiDAR). Only \emph{single-frame}, clustering-driven unsupervised methods apply here.
  \item \textbf{One object class, strong known-radius prior, heavy occlusion / burial}
  $\Rightarrow$ favours \emph{lightweight per-object regressors} over heavy multi-class
  driving-style detectors, and points almost exactly at one paper --- BarrelNet [Yan'24].
\end{enumerate}

\subsection{Supervised (train on synthetic ground truth)}
Grouped by \emph{where each slots into the proposer$+$fit pipeline}:
\begin{itemize}[leftmargin=1.4em,itemsep=3pt,topsep=2pt]
  \item \textbf{(a) Per-object pose / parameter regressor --- a learned ``fit''.}
  \textbf{BarrelNet} [Yan'24]: a PointNet that regresses 6-DOF axis, radius and burial
  fraction from one partial / occluded barrel, trained \emph{purely on synthetic} occluded
  barrels, and beats least-squares fitting under occlusion (cos-sim $0.617\!\to\!0.997$;
  burial-fraction error $25.8\%\!\to\!1.9\%$). This is the closest match to the project's
  problem and drops in as a replacement for the \texttt{normals2step}/LS fit, keeping the
  existing proposer. \emph{No public code} --- reproduce on PointNet++ [Qi'17]. \textbf{Top
  pick (method \#5).}
  \item \textbf{(b) Learned segmentation proposer.} A point-wise net labels barrel
  vs.\ background (or per-instance), then cluster $+$ fit --- replacing the brittle hard-coded
  \texttt{find\_barrel\_clusters} width prior that fails on the cluttered pile. Candidates:
  PointNet++ seg [Qi'17], KPConv [Thomas'19], Point Transformer [Zhao'21], RandLA-Net
  [Hu'20], or the industrial-pipe-specific ResPointNet++ [Yin'21]. Isaac Sim supplies the
  per-point / instance labels for free.
  \item \textbf{(c) End-to-end learned detector (the ``learned LiDAR'' family entry).} Via
  OpenPCDet: PointPillars [Lang'19], PV-RCNN [Shi'20], CenterPoint [Yin'21], or VoteNet
  [Qi'19] for single-object indoor-scale scenes. Heaviest to train and most
  sim-to-real-sensitive; needed for a fair geometric-vs-learned comparison.
  \item \textbf{(d) Primitive-fitting networks (cylinder is a native output).} SPFN [Li'19],
  Point2Cyl [Uy'22] (a barrel \emph{is} an extrusion cylinder; code released), HPNet
  [Yan'21], QuadricsNet [QuadricsNet'23], BPNet [BPNet'23], Point2CAD [Liu'24]. \emph{Caveat
  to test, not assume:} all are trained on clean CAD (ABC / DeepCAD) assuming fairly complete
  surfaces --- the CAD$\to$noisy-single-view-LiDAR$+$occlusion gap is the open question;
  expect to finetune on the synthetic set.
\end{itemize}

\subsection{Unsupervised / self-supervised}
Directly serves the thesis's \emph{annotation / training cost} column:
\begin{itemize}[leftmargin=1.4em,itemsep=3pt,topsep=2pt]
  \item \textbf{(a) Self-supervised pre-train $\to$ few-label finetune (highest novelty).}
  Pre-train on the \emph{unlabeled} raw scans (e.g.\ the 1.23\,M-pt survey scan) with a
  masked-autoencoder --- Occupancy-MAE [Min'23], MAELi [Krispel'24] --- or contrastive
  objective --- PointContrast [Xie'20], DepthContrast [Zhang'21] (single-view depth, a good
  fit for single-frame data) --- then finetune a detector on \emph{very few} labels. Yields a
  \emph{label-efficiency curve} (accuracy vs.\ \# annotated barrels) that geometric methods
  cannot offer.
  \item \textbf{(b) Unsupervised primitive decomposition (no labels, cylinder-native).}
  Recursive deep primitive fitting [Sharma'22], PriFit [PriFit'22], Learnable Earth Parser
  [Loiseau'24] decompose a cloud into primitives via a reconstruction loss; cylinder is a
  primitive. Research-grade but a legitimate unsupervised entry.
  \item \textbf{(c) Single-frame unsupervised detection.} Commonsense-Prototype CPD [Wu'24]
  clusters geometry into proposals and self-trains a detector \emph{without manual boxes} ---
  one of the few that works on a single static frame. (Motion-based LISO / MODEST / OYSTER
  are listed only to record \emph{why} they are excluded here.)
  \item \textbf{(d) Classical unsupervised baseline.} DBSCAN / Euclidean clustering $+$
  region-growing-on-normals --- effectively the project's \emph{current} geometric proposer,
  so it is the unsupervised baseline to beat rather than a new method.
\end{itemize}

\subsection{Recommended ML shortlist (priority order)}
\begin{enumerate}[leftmargin=1.6em,itemsep=2pt,topsep=2pt]
  \item \textbf{BarrelNet-style learned fit} (supervised on synthetic) --- closest prior,
  occlusion-robust, drop-in into proposer$+$fit, cheap to train, directly comparable to the
  four geometric fits. \emph{Do first.}
  \item \textbf{Learned segmentation proposer} (PointNet++ / KPConv / ResPointNet++ on Isaac
  Sim instance labels) --- attacks the exact place geometric methods break (the pile).
  \item \textbf{One OpenPCDet detector} (CenterPoint or PV-RCNN) --- the end-to-end
  learned-LiDAR comparison point.
  \item \textbf{Self-supervised pre-train $+$ few-label finetune} (Occupancy-MAE) --- the
  label-efficiency axis; the most novel contribution beyond a plain bake-off.
  \item \emph{Optional / exploratory:} Point2Cyl transfer test and one unsupervised
  primitive-decomposition entry.
\end{enumerate}

\noindent\textbf{Sequencing caveat.} Every supervised option (1--3) is blocked on the Isaac
Sim data pipeline being up; until then only the unsupervised / pre-trained-transfer paths
(4, plus CAD-pretrained Point2Cyl as a zero-shot test) can touch real data at all.

% =================== 5c. FUTURE PIPELINE ===================
\section{Future pipeline: methods unlocked by Isaac Sim + labelled data}
\label{sec:future}

\emph{(Added 2026-06-30.)} \S\ref{sec:mlsupervision} deliberately \emph{filters down} to what
runs on today's tiny, unlabelled real set. This section is the complementary \emph{expansion}:
once the Isaac Sim pipeline delivers labelled data, all three constraints of
\S\ref{sec:mlsupervision} relax and a much larger menu becomes trainable and \emph{measurable}.

\subsection{What the simulator supplies (and why each unlocks methods)}
Isaac Sim / Replicator gives, per scene, for free: \textbf{(i)} exact 6-DOF pose, radius and
height per barrel; \textbf{(ii)} \emph{true per-barrel occlusion \%} (visible/total surface
points) --- the grounded $x$-axis for the occlusion sweep; \textbf{(iii)} semantic \emph{and}
instance segmentation; \textbf{(iv)} an RTX LiDAR matched to the real sensor and a
\emph{perfectly co-registered} camera; \textbf{(v)} scripted sweeps of barrel count, spacing,
burial depth and mutual occlusion. Each method below is annotated with the label it consumes
and the headline experiment (\S\,Expt-1 occlusion sweep, \S\,Expt-2 calibration-noise) it feeds.

\subsection{(1) Full supervised LiDAR detectors --- now trainable at scale}
The end-to-end detectors that \S\ref{sec:mlsupervision}(c) could only gesture at become the
core ``learned LiDAR'' arm: \textbf{PointPillars} [Lang'19], \textbf{SECOND},
\textbf{PV-RCNN} [Shi'20], \textbf{CenterPoint} [Yin'21], \textbf{VoteNet} [Qi'19], all in one
OpenPCDet harness. \emph{Consumes:} 3D box / instance labels. \emph{Feeds:} Expt-1 (recall +
pose error vs.\ occlusion \%), and the geometric-vs-learned accuracy/cost comparison.

\subsection{(2) Occlusion-specialist learned detection --- the thesis's sharp end}
With true occlusion \% as both a training signal and an evaluation axis, occlusion-aware
methods move from ``interesting'' to \emph{central}:
\begin{itemize}[leftmargin=1.4em,itemsep=2pt,topsep=2pt]
  \item \textbf{BtcDet} [Xu'22] --- predicts occupancy in occluded regions before detecting;
  the sim's occupancy / occlusion GT is exactly its supervision. \emph{The occlusion-question
  workhorse.}
  \item \textbf{Shape completion $\to$ fit} --- PCN [Yuan'18] (and newer completion nets):
  hallucinate the hidden shell, then apply \emph{any} geometric fit --- a learned pre-stage
  that upgrades all four existing methods under occlusion.
  \item \textbf{Amodal detection} --- Isaac Sim can render amodal (whole-object) masks à la
  Amodal-SynthDrive [Sun'24], enabling detectors that regress the \emph{complete} barrel from
  partial evidence. \emph{Feeds:} Expt-1.
\end{itemize}

\subsection{(3) Camera+LiDAR fusion --- now that extrinsics are perfect}
Co-registration makes the whole fusion family (Group~D) trainable, spanning the design axis:
\begin{itemize}[leftmargin=1.4em,itemsep=2pt,topsep=2pt]
  \item \textbf{Camera-first frustum} --- Frustum-PointNets [Qi'18] / -ConvNet [Wang'19]: the
  camera still sees a barrel when its LiDAR return is sparse $\Rightarrow$ the recall-under-
  occlusion hypothesis to \emph{test}, not assume.
  \item \textbf{Decoration / decision-level} --- PointPainting [Vora'20], CLOCs [Pang'20]:
  cheap, modular, plug into OpenPCDet.
  \item \textbf{Deep fusion} --- TransFusion [Bai'22], BEVFusion [Liu'23]: SOTA and, per a
  robustness benchmark [Yu'23], \emph{calibration-robust} (TransFusion loses only
  ${\sim}0.5\%$ mAP at 1\,m misalignment vs.\ $2.3$--$2.9\%$ for hard-projection fusion).
\end{itemize}
\emph{Feeds:} \textbf{Expt-2 (calibration-noise injection)} --- perturb the sim-perfect
extrinsics and measure degradation \emph{per fusion type}; [Yu'23] gives the a-priori ranking
(soft-attention deep fusion $>$ decoration $>$ camera-first frustum) that the thesis can
confirm or overturn on barrels.

\subsection{(4) Sim-to-real domain adaptation --- the bridge you will need}
Because the models above train on \emph{synthetic} clouds, closing the sim-to-real gap is its
own method family (and its own measurable axis --- how much of the gap each technique closes):
\begin{itemize}[leftmargin=1.4em,itemsep=2pt,topsep=2pt]
  \item \textbf{Self-training UDA} --- ST3D [Yang'21] / ST3D++ [Yang'22], STAL3D [Wang'24]:
  pre-train on sim, then pseudo-label + curriculum-retrain on unlabelled real scans; often
  recovers most of the fully-supervised gap.
  \item \textbf{Self-supervised pre-train $\to$ finetune} --- Occupancy-MAE [Min'23], SPOT
  [Yan'23], DepthContrast [Zhang'21]: pre-train on unlabelled real, finetune on sim labels ---
  complements \S\ref{sec:mlsupervision}(a) and yields the label-efficiency curve.
  \item \textbf{Object-level / sensor-model DA} --- match sim LiDAR noise, intensity and
  point density to the real sensor (already partly handled by RTX-LiDAR matching).
\end{itemize}

\subsection{(5) BarrelNet, done properly}
The \S\ref{sec:mlsupervision} top pick graduates: trained on Isaac Sim with \emph{true}
burial-fraction and occlusion-\% supervision (not just synthetic proxies), and validated
against real annotated drums --- the direct, apples-to-apples successor to the geometric fits,
and the closest reproduction of the one published barrel-specific result [Yan'24].

\subsection{Framing note: the cost column flips}
For \emph{sim}-trained methods the manual-annotation cost that the survey keeps as a fairness
column collapses to near-zero; the real cost becomes \emph{simulator build $+$ sim-to-real
adaptation}. The comparison therefore shifts from ``geometric (no training) vs.\ learned
(expensive labels)'' to ``geometric vs.\ learned-on-free-sim-labels, discounted by the
sim-to-real gap'' --- which is precisely why (4) is a first-class part of the future pipeline,
not an afterthought.

% =================== 6. BIBLIOGRAPHY ===================
\section{References (grouped, with links)}

{\small

\subsection*{A. Classical / geometric cylinder detection \& fitting}
\begin{itemize}[leftmargin=1.4em,itemsep=3pt,topsep=2pt]
  \item \textbf{[Rabbani'05]} T.~Rabbani, F.~van den Heuvel. \emph{Efficient Hough
  Transform for Automatic Detection of Cylinders in Point Clouds.} ISPRS Workshop ``Laser
  Scanning 2005''. \,(2-step Gaussian-sphere + 3D Hough --- the method 3DTK
  \texttt{detectCylinder} implements.) \href{https://www.isprs.org/proceedings/XXXVI/3-W19/papers/060.pdf}{[pdf]}
  \item \textbf{[Borrmann'11]} D.~Borrmann, J.~Elseberg, K.~Lingemann, A.~N\"uchter.
  \emph{The 3D Hough Transform for Plane Detection in Point Clouds: A Review and a New
  Accumulator Design.} 3D Research 2(2), 2011. \,(3DTK Hough lineage.)
  \href{https://robotik.informatik.uni-wuerzburg.de/telematics/download/3dresearch2011.pdf}{[pdf]}
  \item \textbf{[Schnabel'07]} R.~Schnabel, R.~Wahl, R.~Klein. \emph{Efficient RANSAC for
  Point-Cloud Shape Detection.} Computer Graphics Forum 26(2):214--226, 2007.
  \href{https://onlinelibrary.wiley.com/doi/10.1111/j.1467-8659.2007.01016.x}{[doi]}
  \item \textbf{[Luk\'acs'98]} G.~Luk\'acs, R.~Martin, D.~Marshall. \emph{Faithful
  Least-Squares Fitting of Spheres, Cylinders, Cones and Tori for Reliable Segmentation.}
  ECCV 1998, LNCS 1406. \href{https://link.springer.com/chapter/10.1007/BFb0055697}{[link]}
  \item \textbf{[Rabbani'06]} T.~Rabbani, F.~van den Heuvel, G.~Vosselman.
  \emph{Segmentation of Point Clouds Using Smoothness Constraint.} ISPRS Commission~V
  Symposium, 2006. \,(Region growing by normal smoothness.)
  \item \textbf{[Nurunnabi'19]} A.~Nurunnabi, Y.~Sadahiro, R.~Lindenbergh, D.~Belton.
  \emph{Robust Cylinder Fitting in Laser Scanning Point Cloud Data.} Measurement
  138:632--651, 2019. \href{https://www.sciencedirect.com/science/article/abs/pii/S0263224119301046}{[link]}
  \item \textbf{[Ara\'ujo'20]} A.~M.~C.~Ara\'ujo, M.~M.~Oliveira. \emph{Connectivity-based
  Cylinder Detection in Unorganized Point Clouds.} Pattern Recognition 100:107161, 2020.
  \href{https://www.sciencedirect.com/science/article/abs/pii/S0031320319304613}{[link]}
  $\cdot$ \href{https://www.inf.ufrgs.br/~oliveira/pubs_files/CD/CD.html}{[code]}
  \item \textbf{[Moradi'21]} S.~Moradi, D.~Laurendeau, C.~Gosselin. \emph{Multiple Cylinder
  Extraction from Organized Point Clouds.} Sensors 21(22):7630, 2021.
  \href{https://www.mdpi.com/1424-8220/21/22/7630}{[link]}
\end{itemize}

\subsection*{B. Learned LiDAR-only detectors}
\begin{itemize}[leftmargin=1.4em,itemsep=3pt,topsep=2pt]
  \item \textbf{[Qi'17]} C.~R.~Qi, L.~Yi, H.~Su, L.~J.~Guibas. \emph{PointNet++: Deep
  Hierarchical Feature Learning on Point Sets in a Metric Space.} NeurIPS 2017.
  \href{https://arxiv.org/abs/1706.02413}{[arXiv:1706.02413]}
  \item \textbf{[Qi'19]} C.~R.~Qi, O.~Litany, K.~He, L.~J.~Guibas. \emph{Deep Hough Voting
  for 3D Object Detection in Point Clouds (VoteNet).} ICCV 2019.
  \href{https://arxiv.org/abs/1904.09664}{[arXiv:1904.09664]}
  \item \textbf{[Lang'19]} A.~H.~Lang, S.~Vora, H.~Caesar, L.~Zhou, J.~Yang, O.~Beijbom.
  \emph{PointPillars: Fast Encoders for Object Detection from Point Clouds.} CVPR 2019.
  \href{https://arxiv.org/abs/1812.05784}{[arXiv:1812.05784]}
  \item \textbf{[Shi'20]} S.~Shi, C.~Guo, L.~Jiang, et al. \emph{PV-RCNN: Point-Voxel
  Feature Set Abstraction for 3D Object Detection.} CVPR 2020.
  \href{https://arxiv.org/abs/1912.13192}{[arXiv:1912.13192]} \,(see also SECOND: Yan et
  al., Sensors 2018.)
  \item \textbf{[Shi'19]} S.~Shi, X.~Wang, H.~Li. \emph{PointRCNN: 3D Object Proposal
  Generation and Detection from Point Cloud.} CVPR 2019.
  \href{https://arxiv.org/abs/1812.04244}{[arXiv:1812.04244]} \,(KITTI val car
  AP$_{3D}$@0.7: Easy 88.88\% / Moderate 78.63\% / Hard 77.38\% --- the occlusion-tier gap
  used as quantitative motivation in \S4.)
  \item \textbf{[Yin'21]} T.~Yin, X.~Zhou, P.~Kr\"ahenb\"uhl. \emph{Center-based 3D Object
  Detection and Tracking (CenterPoint).} CVPR 2021.
  \href{https://arxiv.org/abs/2006.11275}{[arXiv:2006.11275]}
  \item \textbf{[Xu'22]} Q.~Xu, Y.~Zhong, U.~Neumann. \emph{Behind the Curtain: Learning
  Occluded Shapes for 3D Object Detection (BtcDet).} AAAI 2022.
  \href{https://arxiv.org/abs/2112.02205}{[arXiv:2112.02205]}
  \item \textbf{[OpenPCDet]} OpenPCDet --- open-source toolbox implementing SECOND,
  PointPillars, PV-RCNN, CenterPoint, etc. \href{https://github.com/open-mmlab/OpenPCDet}{[github]}
  \item \textbf{[Thomas'19]} H.~Thomas, C.~R.~Qi, J.-E.~Deschaud, B.~Marcotegui,
  F.~Goulette, L.~J.~Guibas. \emph{KPConv: Flexible and Deformable Convolution for Point
  Clouds.} ICCV 2019. \href{https://arxiv.org/abs/1904.08889}{[arXiv:1904.08889]}
  \,(point-convolution segmentation backbone.)
  \item \textbf{[Zhao'21]} H.~Zhao, L.~Jiang, J.~Jia, P.~Torr, V.~Koltun. \emph{Point
  Transformer.} ICCV 2021. \href{https://arxiv.org/abs/2012.09164}{[arXiv:2012.09164]}
  \item \textbf{[Hu'20]} Q.~Hu, B.~Yang, L.~Xie, et al. \emph{RandLA-Net: Efficient Semantic
  Segmentation of Large-Scale Point Clouds.} CVPR 2020.
  \href{https://arxiv.org/abs/1911.11236}{[arXiv:1911.11236]}
  \item \textbf{[Yin'21seg]} C.~Yin, B.~Wang, V.~J.~L.~Gan, M.~Wang, J.~C.~P.~Cheng.
  \emph{Automated Semantic Segmentation of Industrial Point Clouds Using ResPointNet++.}
  Automation in Construction 130:103874, 2021.
  \href{https://www.sciencedirect.com/science/article/abs/pii/S0926580521003253}{[link]}
  $\cdot$ \href{https://github.com/PointCloudYC/ResPointNet2}{[code]} \,(pipes / tanks /
  pumps; the closest industrial-cylinder segmentation precedent.)
\end{itemize}

\subsection*{C. Learned primitive fitting \& shape completion}
\begin{itemize}[leftmargin=1.4em,itemsep=3pt,topsep=2pt]
  \item \textbf{[Li'19]} L.~Li, M.~Sung, A.~Dubrovina, L.~Yi, L.~J.~Guibas.
  \emph{Supervised Fitting of Geometric Primitives to 3D Point Clouds (SPFN).} CVPR 2019.
  \href{https://arxiv.org/abs/1811.08988}{[arXiv:1811.08988]}
  \item \textbf{[Uy'22]} M.~A.~Uy, Y.-Y.~Chang, M.~Sung, P.~Goel, J.~Lambourne, T.~Birdal,
  L.~J.~Guibas. \emph{Point2Cyl: Reverse Engineering 3D Objects from Point Clouds to
  Extrusion Cylinders.} CVPR 2022. \href{https://arxiv.org/abs/2112.09329}{[arXiv:2112.09329]}
  \item \textbf{[Sommer'20]} C.~Sommer, Y.~Sun, E.~Bylow, D.~Cremers. \emph{PrimiTect:
  Fast Continuous Hough Voting for Primitive Detection.} ICRA 2020.
  \href{https://arxiv.org/abs/2005.07457}{[arXiv:2005.07457]}
  \item \textbf{[Yuan'18]} W.~Yuan, T.~Khot, D.~Held, C.~Mertz, M.~Hebert. \emph{PCN: Point
  Completion Network.} 3DV 2018. \href{https://arxiv.org/abs/1808.00671}{[arXiv:1808.00671]}
  \item \textbf{[Sharma'20]} G.~Sharma, D.~Liu, S.~Maji, E.~Kalogerakis, S.~Chaudhuri,
  R.~M\v{e}ch. \emph{ParSeNet: A Parametric Surface Fitting Network for 3D Point Clouds.}
  ECCV 2020. \href{https://arxiv.org/abs/2003.12181}{[arXiv:2003.12181]}
  \item \textbf{[Yan'21]} S.~Yan, Z.~Yang, C.~Ma, H.~Huang, E.~Vouga, Q.~Huang. \emph{HPNet:
  Deep Primitive Segmentation Using Hybrid Representations.} ICCV 2021.
  \href{https://arxiv.org/abs/2105.10620}{[arXiv:2105.10620]}
  $\cdot$ \href{https://github.com/SimingYan/HPNet}{[code]}
  \item \textbf{[QuadricsNet'23]} \emph{QuadricsNet: Learning Concise Representation for
  Geometric Primitives in Point Clouds.} 2023.
  \href{https://arxiv.org/abs/2309.14211}{[arXiv:2309.14211]} \,(lightweight quadric fit;
  handles sparse / noisy data.)
  \item \textbf{[BPNet'23]} \emph{BPNet: B\'ezier Primitive Segmentation on 3D Point Clouds.}
  IJCAI 2023. \href{https://arxiv.org/abs/2307.04013}{[arXiv:2307.04013]}
  \item \textbf{[Liu'24]} \emph{Point2CAD: Reverse Engineering CAD Models from 3D Point
  Clouds.} CVPR 2024. \href{https://arxiv.org/abs/2312.04962}{[arXiv:2312.04962]}
  \,(newer primitive$+$spline reverse-engineering pipeline.)
\end{itemize}

\subsection*{D. Camera--LiDAR fusion}
\begin{itemize}[leftmargin=1.4em,itemsep=3pt,topsep=2pt]
  \item \textbf{[Qi'18]} C.~R.~Qi, W.~Liu, C.~Wu, H.~Su, L.~J.~Guibas. \emph{Frustum
  PointNets for 3D Object Detection from RGB-D Data.} CVPR 2018.
  \href{https://arxiv.org/abs/1711.08488}{[arXiv:1711.08488]}
  \item \textbf{[Wang'19]} Z.~Wang, K.~Jia. \emph{Frustum ConvNet: Sliding Frustums to
  Aggregate Local Point-Wise Features for Amodal 3D Object Detection.} IROS 2019.
  \href{https://arxiv.org/abs/1903.01864}{[arXiv:1903.01864]}
  \item \textbf{[Vora'20]} S.~Vora, A.~H.~Lang, B.~Helou, O.~Beijbom. \emph{PointPainting:
  Sequential Fusion for 3D Object Detection.} CVPR 2020.
  \href{https://arxiv.org/abs/1911.10150}{[arXiv:1911.10150]} \,(see also PointAugmenting,
  Wang et al., CVPR 2021.)
  \item \textbf{[Sindagi'19]} V.~A.~Sindagi, Y.~Zhou, O.~Tuzel. \emph{MVX-Net: Multimodal
  VoxelNet for 3D Object Detection.} ICRA 2019.
  \href{https://arxiv.org/abs/1904.01649}{[arXiv:1904.01649]}
  \item \textbf{[Huang'20]} T.~Huang, Z.~Liu, X.~Chen, X.~Bai. \emph{EPNet: Enhancing Point
  Features with Image Semantics for 3D Object Detection.} ECCV 2020.
  \href{https://arxiv.org/abs/2007.08856}{[arXiv:2007.08856]}
  \item \textbf{[Pang'20]} S.~Pang, D.~Morris, H.~Radha. \emph{CLOCs: Camera-LiDAR Object
  Candidates Fusion for 3D Object Detection.} IROS 2020.
  \href{https://arxiv.org/abs/2009.00784}{[arXiv:2009.00784]}
  \item \textbf{[Bai'22]} X.~Bai, Z.~Hu, X.~Zhu, et al. \emph{TransFusion: Robust
  LiDAR-Camera Fusion for 3D Object Detection with Transformers.} CVPR 2022.
  \href{https://arxiv.org/abs/2203.11496}{[arXiv:2203.11496]}
  \item \textbf{[Liu'23]} Z.~Liu, H.~Tang, A.~Amini, et al. \emph{BEVFusion: Multi-Task
  Multi-Sensor Fusion with Unified Bird's-Eye-View Representation.} ICRA 2023.
  \href{https://arxiv.org/abs/2205.13542}{[arXiv:2205.13542]}
  \item \textbf{[Kumar'25]} S.~Kumar, T.~Brophy, E.~M.~Grua, G.~Sistu, V.~Donzella,
  C.~Eising. \emph{Evaluating the Impact of Weather-Induced Sensor Occlusion on
  BEVFusion for 3D Object Detection.} arXiv, Nov 2025.
  \href{https://arxiv.org/abs/2511.04347}{[arXiv:2511.04347]} \,(nuScenes; controlled
  camera/LiDAR occlusion injection --- quantitative motivation in \S4.)
  \item \textbf{[Wang'25]} H.~Wang, J.~Liu, H.~Dong, Z.~Shao. \emph{A Survey of the
  Multi-Sensor Fusion Object Detection Task in Autonomous Driving.} Sensors 25(9):2794,
  2025. \href{https://www.mdpi.com/1424-8220/25/9/2794}{[link]} \,(qualitative: LiDAR
  limited by occlusion/cost, camera limited by light/weather/depth --- fusion as the
  complementary fix.)
  \item \textbf{[Yu'23]} K.~Yu, T.~Tao, H.~Xie, et al. \emph{Benchmarking the Robustness of
  LiDAR-Camera Fusion for 3D Object Detection.} CVPR-W 2023.
  \href{https://arxiv.org/abs/2205.14951}{[arXiv:2205.14951]} \,(ranks fusion types by
  calibration-noise tolerance --- the a-priori for the Expt-2 calibration-noise injection;
  TransFusion's soft attention degrades least.)
\end{itemize}

\subsection*{E. Occlusion-robust \& barrel / buried-object specific}
\begin{itemize}[leftmargin=1.4em,itemsep=3pt,topsep=2pt]
  \item \textbf{[Yan'24]} J.~Yan, C.~Talegaonkar, N.~Antipa, E.~Terrill, S.~Merrifield.
  \emph{Pose Estimation of Buried Deep-Sea Objects using 3D Vision Deep Learning Models}
  (BarrelNet). OCEANS 2024 --- Halifax. \href{https://arxiv.org/abs/2410.01061}{[arXiv:2410.01061]}
  \,\textbf{--- closest prior work.}
  \item \textbf{[Xu'22]} BtcDet --- ``Behind the Curtain'' (occlusion-aware detection);
  listed under \S B above. \href{https://arxiv.org/abs/2112.02205}{[arXiv:2112.02205]}
  \item \textbf{[Schinagl'22]} D.~Schinagl, G.~Krispel, H.~Possegger, P.~M.~Roth,
  H.~Bischof. \emph{OccAM's Laser: Occlusion-based Attribution Maps for 3D Object
  Detectors on LiDAR Data.} CVPR 2022. \href{https://arxiv.org/abs/2204.06577}{[arXiv:2204.06577]}
  \item \textbf{[Geophys'10]} \emph{Integrated Geophysical Measurements on a Test Site for
  Detection of Buried Steel Drums.} 2010. \,(Non-LiDAR buried-drum context: magnetometry /
  EMI.) \href{https://arxiv.org/abs/1009.5555}{[arXiv:1009.5555]}
  \item \textbf{[Sun'24]} A.~Sun, et al. \emph{Amodal-SynthDrive: A Synthetic Amodal
  Perception Dataset for Autonomous Driving.} 2024.
  \href{https://arxiv.org/abs/2309.06547}{[arXiv:2309.06547]} \,(synthetic amodal / whole-object
  masks --- the template for rendering amodal barrel GT in Isaac Sim.)
\end{itemize}

\subsection*{F. Self-supervised / unsupervised LiDAR representation learning \& detection}
\begin{itemize}[leftmargin=1.4em,itemsep=3pt,topsep=2pt]
  \item \textbf{[Xie'20]} S.~Xie, J.~Gu, D.~Guo, C.~R.~Qi, L.~J.~Guibas, O.~Litany.
  \emph{PointContrast: Unsupervised Pre-training for 3D Point Cloud Understanding.} ECCV
  2020. \href{https://arxiv.org/abs/2007.10985}{[arXiv:2007.10985]}
  \item \textbf{[Zhang'21]} Z.~Zhang, R.~Girdhar, A.~Joulin, I.~Misra. \emph{Self-Supervised
  Pretraining of 3D Features on any Point-Cloud (DepthContrast).} ICCV 2021.
  \href{https://arxiv.org/abs/2101.02691}{[arXiv:2101.02691]}
  $\cdot$ \href{https://github.com/facebookresearch/DepthContrast}{[code]} \,(single-view
  depth --- fits single-frame data.)
  \item \textbf{[Min'23]} C.~Min, et al. \emph{Occupancy-MAE: Self-Supervised Pre-Training
  Large-Scale LiDAR Point Clouds with Masked Occupancy Autoencoders.} IEEE Trans.\
  Intelligent Vehicles, 2023. \href{https://arxiv.org/abs/2206.09900}{[arXiv:2206.09900]}
  \item \textbf{[Krispel'24]} G.~Krispel, et al. \emph{MAELi: Masked Autoencoder for
  Large-Scale LiDAR Point Clouds.} WACV 2024.
  \href{https://arxiv.org/abs/2212.07207}{[arXiv:2212.07207]}
  \item \textbf{[Baur'24]} S.~Baur, et al. \emph{LISO: LiDAR-Only Self-Supervised 3D Object
  Detection.} ECCV 2024. \href{https://arxiv.org/abs/2403.07071}{[arXiv:2403.07071]}
  \,(\emph{needs sequential / motion} cues --- recorded here as \emph{not} applicable to the
  single static scan.)
  \item \textbf{[Wu'24]} H.~Wu, et al. \emph{Commonsense Prototype for Outdoor Unsupervised
  3D Object Detection (CPD).} CVPR 2024. \href{https://arxiv.org/abs/2404.16493}{[arXiv:2404.16493]}
  \,(single-frame, clustering-driven self-training.)
\end{itemize}

\subsection*{G. Unsupervised primitive decomposition (no labels)}
\begin{itemize}[leftmargin=1.4em,itemsep=3pt,topsep=2pt]
  \item \textbf{[Sharma'22]} \emph{Unsupervised Recursive Deep Fitting of 3D Primitives to
  Points.} Computers \& Graphics, 2022.
  \href{https://www.sciencedirect.com/science/article/abs/pii/S0097849321002314}{[link]}
  \item \textbf{[PriFit'22]} \emph{PriFit: Learning to Fit Primitives Improves Few-Shot Point
  Cloud Segmentation.} Eurographics / SGP 2022.
  \href{https://arxiv.org/abs/2112.13942}{[arXiv:2112.13942]}
  \item \textbf{[Loiseau'24]} R.~Loiseau, E.~Vincent, M.~Aubry, L.~Landrieu. \emph{Learnable
  Earth Parser: Discovering 3D Prototypes in Aerial Scans.} CVPR 2024.
  \href{https://arxiv.org/abs/2304.09704}{[arXiv:2304.09704]} \,(unsupervised reconstruction
  $\to$ primitive prototypes.)
\end{itemize}

\subsection*{H. Synthetic data, sim-to-real \& domain adaptation for learned LiDAR detection}
\begin{itemize}[leftmargin=1.4em,itemsep=3pt,topsep=2pt]
  \item \textbf{[Saleh'20]} K.~Saleh, et al. \emph{Manual-Label Free 3D Detection via an
  Open-Source Simulator.} 2020. \href{https://arxiv.org/abs/2011.07784}{[arXiv:2011.07784]}
  \,(simulator-trained detector with LiDAR-guided sampling; supports the Isaac-Sim
  train$\to$real-test plan.)
  \item \textbf{[Yang'21]} J.~Yang, S.~Shi, Z.~Wang, H.~Li, X.~Qi. \emph{ST3D: Self-training
  for Unsupervised Domain Adaptation on 3D Object Detection.} CVPR 2021.
  \href{https://arxiv.org/abs/2103.05346}{[arXiv:2103.05346]}
  $\cdot$ \href{https://github.com/CVMI-Lab/ST3D}{[code]} \,(canonical sim/real-to-real UDA;
  pseudo-label memory + curriculum augmentation.)
  \item \textbf{[Yang'22]} J.~Yang, S.~Shi, Z.~Wang, H.~Li, X.~Qi. \emph{ST3D++: Denoised
  Self-training for Unsupervised Domain Adaptation on 3D Object Detection.} IEEE T-PAMI 2022.
  \href{https://arxiv.org/abs/2108.06682}{[arXiv:2108.06682]}
  \item \textbf{[Wang'24]} Y.~Wang, et al. \emph{STAL3D: Unsupervised Domain Adaptation for 3D
  Object Detection via Collaborating Self-Training and Adversarial Learning.} 2024.
  \href{https://arxiv.org/abs/2406.19362}{[arXiv:2406.19362]}
  \item \textbf{[Yan'23]} X.~Yan, et al. \emph{SPOT: Scalable 3D Pre-training via Occupancy
  Prediction for Learning Transferable 3D Representations.} 2023.
  \href{https://arxiv.org/abs/2309.10527}{[arXiv:2309.10527]} \,(occupancy-based pre-training;
  bridges sim label $\to$ real finetune.)
\end{itemize}

} % end \small

\vspace{6pt}
\hrule
\vspace{4pt}
{\footnotesize\emph{Citation note.} References were compiled via literature search on
2026-05-30. Authors, venues, and arXiv identifiers were checked against the linked
sources, but please re-verify each entry against its link before quoting it in the final
thesis. This file lives in \texttt{researchwrite/} and is regenerated as the survey
evolves; the companion index is the \texttt{researchwrite-materials} memory.}

% =================== 7. IMPLEMENTATIONS ===================
\clearpage
\section{Off-the-shelf implementations (GitHub) --- what to reuse}

The plan is to \textbf{reuse code, not reimplement}. For each method (keyed to the table
in \S3) the table below gives the \emph{recommended} repository --- the most-starred,
actively maintained option --- with lighter or alternative repos in \emph{Notes}. Star
counts (\ding{72}), language and licence are as of 2026-05-30.

\textbf{Two caveats.} (i)~\textbf{BarrelNet} [Yan'24], the closest prior work, has
\emph{no public code release} --- it must be reproduced on top of a PointNet++ repo.
(ii)~A few official repos are \emph{archived} (read-only on GitHub) but still compile and
run.

\vspace{4pt}
\noindent\colorbox{cH!55}{%
\begin{minipage}{\dimexpr\textwidth-2\fboxsep\relax}
\textbf{Toolchain shortcut.}\; Standardise on
\href{https://github.com/open-mmlab/OpenPCDet}{open-mmlab/OpenPCDet}. PointPillars,
SECOND, PV-RCNN and CenterPoint all ship inside it; BtcDet is built on it; and CLOCs and
the recommended PointPainting repo plug into it. Adopting one framework puts almost all
the learned-LiDAR methods plus several fusion methods in a \emph{single} codebase --- little
to no from-scratch work, one training/eval harness, and an easy bridge to the project's
\texttt{predictions.json} contract.
\end{minipage}}

\begin{landscape}
\footnotesize
\setlength{\extrarowheight}{2.5pt}
\renewcommand{\arraystretch}{1.05}

\begin{longtable}{@{}R{3.5cm} R{5.0cm} R{1.2cm} R{2.7cm} R{9.6cm}@{}}
\rowcolor{gray!35}
\textbf{Method (cf.\ \S3)} & \textbf{Recommended GitHub repo} & \textbf{\ding{72} Stars} & \textbf{Lang / Licence} & \textbf{Notes (alternatives \& caveats)}\\
\endhead

\rowcolor{black!8}\multicolumn{5}{@{}l}{\textbf{A.\quad LiDAR-only --- classical / geometric}}\\
Euclidean clustering + fit &
\href{https://github.com/PointCloudLibrary/pcl}{PointCloudLibrary/pcl} & 11.0k & C++ / BSD &
\texttt{EuclideanClusterExtraction} + \texttt{SACSegmentation} (cylinder model) in one library. Python route: \href{https://github.com/isl-org/Open3D}{isl-org/Open3D} (13.6k, MIT) for clustering, paired with a cylinder fit.\\

Region growing + curvature &
\href{https://github.com/PointCloudLibrary/pcl}{PointCloudLibrary/pcl} & 11.0k & C++ / BSD &
\texttt{pcl::RegionGrowing} (smoothness / curvature constraint) --- the [Rabbani'06] method.\\

Randomized 2-step Hough --- 3DTK \emph{(baseline)} &
\href{https://github.com/JMUWRobotics/3DTK}{JMUWRobotics/3DTK} & 131 & C++ &
\textbf{Already vendored.} Maintained mirror; provides \texttt{detectCylinder}.\\

RANSAC cylinder fit &
\href{https://github.com/leomariga/pyRANSAC-3D}{leomariga/pyRANSAC-3D} & 653 & Python / Apache &
Pure-Python \texttt{Cylinder()} --- easiest pull for this stack. C++ reference: PCL \texttt{SACMODEL\_CYLINDER} (11.0k).\\

Efficient RANSAC (multi-primitive) &
\href{https://github.com/CGAL/cgal}{CGAL/cgal} & 5.9k & C++ / GPL+comm. &
``Point-Set Shape Detection'' = the [Schnabel'07] algorithm. Lighter standalone: \href{https://github.com/LiangliangNan/Easy3D}{LiangliangNan/Easy3D} (1.6k, GPL).\\

Robust least-squares cylinder fit &
\href{https://github.com/xingjiepan/cylinder_fitting}{xingjiepan/cylinder\_fitting} & 53 & Python / BSD &
pip-installable least-squares cylinder fit (Eberly). \emph{Fit stage only} --- needs a proposer.\\

Connectivity / projection detection &
\href{https://github.com/abner-math/CylinderDetection}{abner-math/CylinderDetection} & 39 & C++ / GPL &
Author's official code for [Ara\'ujo'20].\\

\rowcolor{black!8}\multicolumn{5}{@{}l}{\textbf{B.\quad LiDAR-only --- learned}}\\
PointNet++ \,(basis for BarrelNet) &
\href{https://github.com/yanx27/Pointnet_Pointnet2_pytorch}{yanx27/Pointnet\_Pointnet2\_pytorch} & 4.9k & Python / MIT &
Most-popular PyTorch PointNet++. \textbf{BarrelNet [Yan'24] has no public code --- reproduce here.} Alt: \href{https://github.com/erikwijmans/Pointnet2_PyTorch}{erikwijmans/Pointnet2\_PyTorch} (1.8k).\\

VoteNet &
\href{https://github.com/facebookresearch/votenet}{facebookresearch/votenet} & 1.8k & Python / MIT &
Official deep-Hough-voting detector. Repo \emph{archived} (read-only) but runs.\\

PointPillars / SECOND / PV-RCNN &
\href{https://github.com/open-mmlab/OpenPCDet}{open-mmlab/OpenPCDet} & 5.6k & Python / Apache &
All three (+ Part-A2, etc.) in one maintained hub. Alt: \href{https://github.com/open-mmlab/mmdetection3d}{open-mmlab/mmdetection3d} (6.4k).\\

CenterPoint &
\href{https://github.com/tianweiy/CenterPoint}{tianweiy/CenterPoint} & 2.2k & Python / MIT &
Official center-based detector; also bundled inside OpenPCDet.\\

BtcDet --- occlusion-aware &
\href{https://github.com/Xharlie/BtcDet}{Xharlie/BtcDet} & 201 & Python / Apache &
Official; built on OpenPCDet 0.3.0 --- drops into the same toolchain. \textbf{Best occlusion specialist.}\\

Learned primitive fitting &
\href{https://github.com/lingxiaoli94/SPFN}{lingxiaoli94/SPFN} & 184 & Python / MIT &
SPFN (multi-primitive). Also \href{https://github.com/mikacuy/point2cyl}{mikacuy/point2cyl} (65, extrusion cylinders) and \href{https://github.com/c-sommer/primitect}{c-sommer/primitect} (33, C++).\\

Shape completion (PCN, pre-stage) &
\href{https://github.com/wentaoyuan/pcn}{wentaoyuan/pcn} & 482 & Python / MIT &
Official (TensorFlow). PyTorch reimpl: \href{https://github.com/qinglew/PCN-PyTorch}{qinglew/PCN-PyTorch} (212).\\

\rowcolor{black!8}\multicolumn{5}{@{}l}{\textbf{C.\quad Camera + LiDAR fusion}}\\
Camera-first frustum &
\href{https://github.com/charlesq34/frustum-pointnets}{charlesq34/frustum-pointnets} & 1.7k & Python / Apache &
Official Frustum-PointNets (TF). F-ConvNet: \href{https://github.com/Gorilla-Lab-SCUT/frustum-convnet}{Gorilla-Lab-SCUT/frustum-convnet} (247, PyTorch).\\

LiDAR-first $\to$ image verification &
\href{https://github.com/pangsu0613/CLOCs}{pangsu0613/CLOCs} & 405 & Python / MIT &
Decision-level fusion on top of SECOND / PV-RCNN (OpenPCDet-compatible).\\

Point decoration (PointPainting) &
\href{https://github.com/Song-Jingyu/PointPainting}{Song-Jingyu/PointPainting} & 314 & Python / MIT &
mmsegmentation + OpenPCDet, maintained. Alt: \href{https://github.com/AmrElsersy/PointPainting}{AmrElsersy/PointPainting} (170, real-time BiSeNetv2).\\

Deep multimodal fusion &
\href{https://github.com/mit-han-lab/bevfusion}{mit-han-lab/bevfusion} & 3.2k & Python / Apache &
Most-popular fusion repo; \emph{archived} but runs. Alt: \href{https://github.com/XuyangBai/TransFusion}{XuyangBai/TransFusion} (760); MVX-Net via mmdetection3d (6.4k); \href{https://github.com/happinesslz/EPNet}{happinesslz/EPNet} (254).\\

\end{longtable}
\end{landscape}

\noindent{\footnotesize Star counts are a popularity snapshot (GitHub API, 2026-05-30) and
drift over time; ``archived'' repos are read-only but still clone/build. Where a method has
several implementations, the most-starred actively-maintained one is recommended and the
rest are listed as alternatives.}

\end{document}
